% ===== PRÉAMBULE CANONIQUE — à coller VERBATIM en tête de CHAQUE .tex =====
\documentclass[11pt,a4paper]{article}
\usepackage[utf8]{inputenc}\usepackage[T1]{fontenc}\usepackage{lmodern}
\usepackage[french]{babel}
\usepackage{amsmath,amssymb,mathtools}\usepackage{icomma}
\usepackage[version=4]{mhchem}          % \ce{...} : équations & formules
\usepackage{siunitx}\sisetup{locale=FR} % \qty{}{}, \si{}, \ang{}
\usepackage{chemfig}                    % \chemfig{...} : structures développées
\usepackage{xcolor}\usepackage{colortbl}
\usepackage{tikz}\usepackage{pgfplots}\pgfplotsset{compat=1.18}
\usepackage[most]{tcolorbox}\usepackage{enumitem}\usepackage{array,booktabs}
\usepackage[a4paper,margin=2.2cm]{geometry}
\usetikzlibrary{arrows.meta,calc,angles,quotes,positioning,patterns,backgrounds,shapes.geometric}
\definecolor{primary}{RGB}{222,49,99}\definecolor{primarydark}{RGB}{150,22,55}
\definecolor{defcol}{RGB}{0,114,178}\definecolor{methcol}{RGB}{52,92,148}
\definecolor{excol}{RGB}{0,135,90}\definecolor{attncol}{RGB}{200,40,55}
\tcbset{boxbase/.style={enhanced,breakable,boxrule=0.9pt,arc=2.5pt,
  left=3mm,right=3mm,top=2.3mm,bottom=2.3mm,fonttitle=\bfseries,coltitle=white}}
% --- boîtes TUTORIEL (noms EXACTS, ne pas renommer) ---
\newtcolorbox{cle}[1][]{boxbase,colback=defcol!6,colframe=defcol,
  title={\textsc{À retenir}\ifx\relax#1\relax\relax\else\textmd{ : \itshape #1}\fi}}
\newtcolorbox{methode}[1][]{boxbase,colback=methcol!7,colframe=methcol,sharp corners,
  title={\textsc{Méthode}\ifx\relax#1\relax\relax\else\textmd{ --- \itshape #1}\fi}}
\newtcolorbox{exemplebox}[1][]{boxbase,colback=excol!5,colframe=excol,
  title={\textsc{Exemple}\ifx\relax#1\relax\relax\else\textmd{ : \itshape #1}\fi}}
\newtcolorbox{piege}[1][]{boxbase,colback=attncol!6,colframe=attncol,
  title={\textsc{Piège}\ifx\relax#1\relax\relax\else\textmd{ : \itshape #1}\fi}}
\newtcolorbox{remarque}[1][]{boxbase,colback=black!4,colframe=black!45,
  title={\textsc{Remarque}\ifx\relax#1\relax\relax\else\textmd{ : \itshape #1}\fi}}
% --- boîtes EXERCICE / CORRIGÉ (noms EXACTS, auto-comptées) ---
\newtcolorbox[auto counter]{exo}[1][]{boxbase,colback=primary!4,colframe=primary,
  title={\textsc{Exercice}~\thetcbcounter\ifx\relax#1\relax\relax\else\textmd{ --- \itshape #1}\fi}}
\newtcolorbox[auto counter]{corrige}[1][]{boxbase,colback=excol!4,colframe=excol,
  title={\textsc{Corrigé}~\thetcbcounter\ifx\relax#1\relax\relax\else\textmd{ --- \itshape #1}\fi}}
\newcommand{\R}{\mathbb{R}}\newcommand{\N}{\mathbb{N}}\newcommand{\Z}{\mathbb{Z}}\newcommand{\ds}{\displaystyle}
% (un \usepackage{fancyhdr}+en-tête est possible ; il est ignoré côté web)
% =========================================================================
\begin{document}
\newcommand{\AX}[2]{\ensuremath{\mathrm{AX_{#1}E_{#2}}}}
\begin{corrige}[la famille bipyramide trigonale]
Les doublets libres se logent en position équatoriale ; on masque les E pour lire la forme.
\begin{enumerate}[label=\alph*)]
  \item \ce{SF4} : \AX{4}{1} $\Rightarrow$ \textbf{à bascule} (papillon).
  \item \ce{ClF3} : \AX{3}{2} $\Rightarrow$ \textbf{en T}.
  \item \ce{XeF2} : \AX{2}{3} $\Rightarrow$ \textbf{linéaire}.
\end{enumerate}
\end{corrige}

\begin{corrige}[la famille octaédrique]
\begin{enumerate}[label=\alph*)]
  \item \ce{SF6} : \AX{6}{0} $\Rightarrow$ \textbf{octaédrique}.
  \item \ce{BrF5} : \AX{5}{1} $\Rightarrow$ \textbf{pyramide à base carrée}.
  \item \ce{XeF4} : \AX{4}{2} $\Rightarrow$ \textbf{plane carrée} (les $2$ doublets se placent en
        positions opposées, à $180^\circ$).
\end{enumerate}
\end{corrige}

\begin{corrige}[pourquoi le doublet va-t-il à l'équateur ?]
\begin{enumerate}[label=\alph*)]
  \item En position \textbf{axiale} : $3$ voisins à $90^\circ$. En position \textbf{équatoriale} :
        seulement $2$ voisins à $90^\circ$.
  \item Le doublet libre, qui repousse le plus fort, choisit la position de \emph{moindre}
        répulsion, donc l'\textbf{équatoriale}. Il reste alors $2$ liaisons axiales et $2$
        équatoriales : la molécule prend la forme \textbf{à bascule} (\ce{SF4}).
\end{enumerate}
\end{corrige}

\begin{corrige}[la fermeture progressive de l'angle]
\begin{enumerate}[label=\alph*)]
  \item \ce{CH4} : \AX{4}{0}, $109{,}5^\circ$ ; \ce{NH3} : \AX{3}{1}, $107^\circ$ ; \ce{H2O} :
        \AX{2}{2}, $104{,}5^\circ$.
  \item Les répulsions décroissent dans l'ordre (non liant--non liant) $>$ (non liant--liant) $>$
        (liant--liant). Chaque doublet libre repousse donc plus fort qu'une liaison et
        \textbf{comprime} les angles : $0$ doublet libre (\ce{CH4}) $\to$ angle intact ;
        $1$ (\ce{NH3}) $\to$ angle réduit à $107^\circ$ ; $2$ (\ce{H2O}) $\to$ réduit à
        $104{,}5^\circ$.
\end{enumerate}
\end{corrige}

\begin{corrige}[deux molécules linéaires… pour deux raisons opposées]
\begin{enumerate}[label=\alph*)]
  \item \ce{BeCl2} : \AX{2}{0} $\Rightarrow$ \textbf{linéaire}.
  \item \ce{SnCl2} : \AX{2}{1} $\Rightarrow$ \textbf{coudée}.
  \item \ce{ICl2-} : \AX{2}{3} $\Rightarrow$ \textbf{linéaire}.
  \item \ce{PCl6-} : \AX{6}{0} $\Rightarrow$ \textbf{octaédrique}.
  \item \ce{BeCl2} (\AX{2}{0}) est linéaire car il n'a \emph{aucun} doublet libre : ses $2$
        liaisons s'écartent à $180^\circ$. \ce{ICl2-} (\AX{2}{3}) est linéaire pour une raison
        \emph{opposée} : il a $5$ sites (bipyramide trigonale), et ses $3$ doublets libres occupent
        les $3$ positions \emph{équatoriales} ; il ne reste que les $2$ liaisons \emph{axiales}, à
        $180^\circ$. Même forme finale, mécanisme différent.
\end{enumerate}
\end{corrige}

\end{document}
